\section{Discussion}

In this section, we discuss the system's generality, limitations, and directions for future improvements.

% \subsection{Lessons Learned}

% During the experiments and user evaluations, we collected failure cases of LLM at various stages of graph construction and analyzed their causes. We identified several common issues and their reasons when LLM constructs graphs:

% \textbf{Incomplete entity extraction.} In our experiments, we found that entity information often comes from multiple articles. When entities fail to be successfully extracted from certain articles, it results in the loss of partial semantic information of these entities, thus affecting the semantic completeness of the graph. This issue prompts us to consider how to enhance the comprehensiveness of entity extraction to ensure the integrity of the graph.

% \textbf{Missing relationships.} Relationships typically originate from specific paragraphs in articles. When these paragraphs fail to extract relationships successfully, it leads to the absence of direct relationships between two entities, affecting the relational completeness of the graph. We recognize the need to develop more robust relationship extraction methods to mitigate this issue.

% \textbf{Incorrect relationships or entities.} Due to the shortcomings of LLM in text understanding or hallucinations, entity or relationship identification errors may occur. For example, misidentifying Entity A as Entity B. This suggests that we need to improve the accuracy and robustness of the model to reduce misidentification.

% \textbf{Duplicate entities.} During entity extraction, due to the LLM's insufficient understanding of names, proper nouns, or hallucinations, different descriptions of the same entity may be treated as two separate entities (e.g., "Donald Trump" and "Trump", "Sam Altman" and "Altman"). These entities carry parts of local relationships and global semantics but are incomplete. This issue prompts us to explore more effective deduplication strategies to enhance the precision of the graph.

% \subsection{System Generalizability}

\textbf{System Generalizability:} Although we adopted Microsoft's \textit{GraphRAG} \cite{edge2024local} as an integration tool in our study, all designs are based on our proposed general architecture. This architecture is compatible with other mainstream GraphRAG systems. For instance, the six stages of our architecture are reflected in all systems, where recall, entity, and relationship are present and serve as the core functionalities. While topic reports may have different meanings and expressions in different systems, we regard them as a type of global summary recall. Other systems may have additional types of recall, leading to extra stages, but this does not conflict with the GraphRAG architecture.

% Our system is not limited to analyzing question-answering on text datasets and can be applied to various scenarios. For example, it can assist in observing graph-based inference processes and play a role in tasks such as knowledge discovery and relationship mining. It is easy to deploy our system on a local machine, which only requires an LLM (e.g. GPT-4) that supports JSON formatted output, as well as a text embedding model that captures text similarity. Additionally, users can provide their own corpus for graph construction and analysis.

% \subsection{Limitations and Future Work}

% \textbf{Limitations and Future Work} Despite the satisfactory performance of our system, there are still some limitations and potential research directions for the future. Here, we list two directions we deem important.

\textbf{Enhancing temporal analysis.} While our system demonstrates satisfactory performance, dynamic changes in entity relationships or topics within a graph, such as the relationship between A and B evolving from friends to enemies, require further investigation. Therefore, it is essential to develop advanced temporal analysis tools to better guide users in exploring the temporal evolution of graphs.

\textbf{Expanding recall capabilities.} Although we currently support entity, relationship, and topic views in mainstream GraphRAG systems with promising results, the development of future GraphRAG systems will necessitate supporting more types of recall to accommodate a wider range of scenarios. This expansion represents another critical direction for improving system capabilities.

%%%%%%%%%%%%%%%%%%%%%%

% 8 讨论

% 在本节中，我们总结了从后续实验中发现的一些常见图质量问题及其原因，并探讨了系统的通用性、局限性以及未来改进的方向。

% 8.1 学到的经验

% 增强对时间变化信息的分析能力。我们认识到图谱中可能存在动态变化的实体关系或话题，例如A和B的关系从好友演变为敌人。因此，我们需要开发先进的时效性分析工具，以便更好地引导用户探索图谱的时序演化。

% 支持更多类型的召回。我们对当前主流 GraphRAG 系统支持的实体、关系和话题视图进行了支持，但是随着未来 GraphRAG 系统的发展，我们将会支持更多类型的召回，以支持更丰富的场景。

% 8.2 通用性

% 我们的工作基于GraphRAG框架，旨在利用该框架实现我们提出的通用GraphRAG架构。尽管我们在研究中采用了GraphRAG作为集成工具，但我们并未对其进行特殊定制，所有设计均基于我们提出的通用架构。该架构与其他主流GraphRAG系统如OpenSPG-KAG、LightRAG等兼容。例如，我们架构的六个阶段在所有系统中均有体现，其中的召回（Recall）、实体（Entity）、关系（Relationship）在这些系统中都存在且作为功能的核心。虽然话题报告（Topic Report）在不同系统中可能具有不同的意义和表达，但我们将其视为一种全局总结类型的召回。在其他系统中可能存在其他类型的召回，因此可能存在额外的阶段，但这并不与GraphRAG架构冲突。

% 我们的系统不仅限于新闻数据集问答的分析，还可以应用于多种场景。例如，它可以辅助观察基于图的推理过程，以及在复杂数据集中的知识发现和关系挖掘等任务中发挥作用。

% 8.3 局限性

% 尽管我们系统的表现令人满意，但仍然存在一些局限性，以及未来可能的研究方向。

% 开放式探索引导不足。我们的系统主要支持用户进行目标导向性的探索，即用户带着明确的问题和目标进行分析。然而，对于无目的的开放式探索，虽然系统提供了全局话题和局部实体关系的探索功能，但缺乏对用户的引导和提示。我们需要增强系统的引导能力，以便更好地支持开放式探索。

% 时效性分析能力不足。我们的系统专注于分析静态信息构成的语料，对于那些存在时间变化的语料，其图谱中可能存在具有时效性的实体关系或话题（例如，A和B一开始是好友，但随着时间推移可能成为敌人）。当前系统无法很好地引导用户进行图谱时序演化的探索，这提示我们需要开发时效性分析工具。

% 复杂推理分析能力不足。对于需要基于图上多跳、多步骤或多智能体推理的复杂推理流程中图质量问题的分析，我们的系统目前无法很好地支持。我们需要提升系统的复杂推理能力，以便更好地处理复杂的图谱分析任务。

% 用户交互界面复杂性。尽管我们的系统提供了丰富的可视化工具，并尽可能简化了系统设计，但由于GraphRAG系统本身的复杂性，对于非技术用户而言，界面的复杂性可能增加学习成本，使得他们难以充分利用系统的功能进行探索和分析。这提示我们需要优化用户界面，提高其易用性和友好性，以降低用户的学习门槛。